% -----------------------------------------------------------
\section{Quicken, Quick}
% -----------------------------------------------------------
	\label{QUICKEN}
	
% % ----------------------
% \subsection{See also}
% % ----------------------

% ----------------------
\subsection{1828 American Dictionary of the English Language} % *1828
% ----------------------
	\label{QUICKEN-1828}
	
	% -----
	\subsubsection{Quicken (transitive)}
	% -----

		\begin{compactitem}
			\item[1] Primarily, to make alive; to vivify; to revive or resuscitate, as from death or an inanimate state.
			\item[2] To make alive in a spiritual sense; to communicate a principle of grace to.
			\item[3] To hasten; to accelerate; as, to \textit{quicken} motion, speed or flight.
			\item[4] To sharpen; to give keener perception to; to stimulate; to incite; as, to \textit{quicken} the appetite or taste; to \textit{quicken} desires.
			\item[5] To revive; to cheer; to reinvigorate; to refresh by new supplies of comfort or grace.
		\end{compactitem}
		
	% -----
	\subsubsection{Quicken (intransitive)}
	% -----
	
		\begin{compactitem}
			\item[1] To become alive.
			\item[2] To move with rapidity or activity.
		\end{compactitem}
		
	% -----
	\subsubsection{Quick (adjective)}
	% -----
	
		\begin{compactitem}
			\item[1] Primarily, alive; living; opposed to dead or unanimated; as \textit{quick} flesh.
			\item[2] Swift; hasty; done with celerity; as \textit{quick} dispatch.
			\item[3] Speedy; done or occurring in a short time; as a \textit{quick} return of profits.
			\item[4] Active; brisk; nimble; prompt ready. He is remarkably \textit{quick} in his motions. He is a man of \textit{quick} parts.
			\item[5] Moving with rapidity or celerity; as \textit{quick} time in music.
		\end{compactitem}

% % ----------------------
% \subsection{Oxford English Dictionary} % *OED
% % ----------------------
	% \label{QUICKEN-OED}

	% \begin{compactitem}
		% \item[]
	% \end{compactitem}
	
% % ----------------------
% \subsection{Buck's Theological Dictionary (1833)} % *BTD
% % ----------------------
	% \label{QUICKEN-BTD}

	% \begin{compactdesc}
		% \item[PAGE\#]
	% \end{compactdesc}
	
% % ----------------------
% \subsection{Restoration Lexicon} % *RL *RESTORATION LEXICON
% % ----------------------
	% \label{QUICKEN-OED}

	% \begin{compactitem}
		% \item[]
	% \end{compactitem}

% ----------------------
\subsection{Scriptural Usage} % *SCRIPTURES
% ----------------------
	\label{QUICKEN-SCRIPTURES}

	% % -----
	% \subsubsection{Old Testament} % *OT
	% % -----
		% \label{QUICKEN-OT}
	
		% % --
		% \paragraph{KJV}
		% % --
			% \label{QUICKEN-OT-KJV}
	
			% \begin{tabular}{ll}
				% Total occurrences in the OT & 
			% \end{tabular}
			
			% \begin{compactdesc}
				% \item[]
			% \end{compactdesc}
			
		% % --
		% \paragraph{Hebrew Bible}
		% % --
			% \label{QUICKEN-HEBREW}
		
			% %
			% \subparagraph{\texorpdfstring{\he{sample word}}{}} % paste actual hebrew text here
			% %
				% \label{PAGA} % whatever is in step bible without periods
			
				% \begin{compactdesc}
					% \item[HALOT , PDF ] \textbf{} % Use the 1994 PDF
						% \begin{compactitem}
							% \item[1]
						% \end{compactitem}
					% \item[BDB , PDF ] \textbf{}
						% \begin{compactitem}
							% \item[1]
						% \end{compactitem}
					% \item[DCH PDF ] \textbf{} % don't include the actual page number, they repeat from volume to volume
						% \begin{compactitem}
							% \item[1]
						% \end{compactitem}
				% \end{compactdesc}
				
				% \begin{compactdesc}
					% \item[Some OT uses] \textbf{}
						% \begin{compactdesc}
							% \item[]
						% \end{compactdesc}
				% \end{compactdesc}
			
		% % --
		% \paragraph{Septuagint}
		% % --
			% \label{QUICKEN-LXX}
		
			% %
			% \subparagraph{\texorpdfstring{\gr{sample word}}{}}
			% %
		
	% -----
	\subsubsection{New Testament} % *NT
	% -----
		\label{QUICKEN-NT}
	
		% --
		\paragraph{KJV}
		% --
			\label{QUICKEN-NT-KJV}		
	
			% \begin{tabular}{ll}
				% Total occurrences in the NT & 
			% \end{tabular}
			
			\begin{compactdesc}
				\item[{\hyperref[JOHN-5-21]{John 5:21}}]
			\end{compactdesc}
			
		% --
		\paragraph{Greek New Testament} % *GREEK
		% --
			\label{QUICKEN-GREEK}
		
			%
			\subparagraph{\texorpdfstring{\gr{ζωοποιέω}}{}}
			%
				\label{ZOOPOIEO}
				
				% \begin{tabular}{ll}
					% Total occurrences in the NT & 
				% \end{tabular}
				
				\begin{compactdesc}
					\item[Some NT uses] \textbf{} % Change to "all" if they're all listed
						\begin{compactdesc}
							\item[{\hyperref[JOHN-5-21]{John 5:21}}]
						\end{compactdesc}
				\end{compactdesc}
			
				\begin{compactdesc}
				
					\item[BDAG 381b] \textbf{}
						\begin{compactitem}
							\item[1] \textbf{to cause to live, \textit{make alive, give life to}}, especially in a transcendent sense 
							\item[1.a] subject God, Christ, and Spirit...
								\begin{compactitem}
									\item Many examples follow in the NT.
								\end{compactitem}
							\item[1.b] of things, circumstances, words...
							\item[2] \textbf{to keep alive, \textit{sustain life}}
						\end{compactitem}
						
					% \item[Brill , PDF ] \textbf{}
						% \begin{compactitem}
							% \item 
						% \end{compactitem}
						
					% \item[CGL , PDF ] \textbf{}
						% \begin{compactitem}
							% \item 
						% \end{compactitem}
						
					% \item[LSJ , PDF ] \textbf{}
						% \begin{compactitem}
							% \item 
						% \end{compactitem}
						
					% \item[TDNT ] \textbf{}
						% \begin{compactitem}
							% \item 
						% \end{compactitem}
						
					% \item[TDNTA ] \textbf{}
						% \begin{compactitem}
							% \item 
						% \end{compactitem}
						
					% \item[NIDNTT ] \textbf{}
						% \begin{compactitem}
							% \item 
						% \end{compactitem}
						
					% \item[NIDNTTE ] \textbf{}
						% \begin{compactitem}
							% \item 
						% \end{compactitem}
						
				\end{compactdesc}
		
	% % -----
	% \subsubsection{Book of Mormon} % *BOM
	% % -----
		% \label{QUICKEN-BOM}
	
		% \begin{tabular}{ll}
			% Total occurrences in the BOM & 
		% \end{tabular}
		
		% \begin{compactdesc}
			% \item[]
		% \end{compactdesc}
		
	% % -----
	% \subsubsection{Doctrine and Covenants} % *DC
	% % -----
		% \label{QUICKEN-DC}
	
		% \begin{tabular}{ll}
			% Total occurrences in the DC & 
		% \end{tabular}
		
		% \begin{compactdesc}
			% \item[]
		% \end{compactdesc}
		
	% % -----
	% \subsubsection{Pearl of Great Price} % *PGP
	% % -----
		% \label{QUICKEN-PGP}
	
		% \begin{tabular}{ll}
			% Total occurrences in the PGP & 
		% \end{tabular}
		
		% \begin{compactdesc}
			% \item[]
		% \end{compactdesc}
		
% % ----------------------
% \subsection{Key Phrases} % *KEYPHRASES *PHRASES
% % ----------------------
	% \label{QUICKEN-PHRASES}

	% \begin{compactdesc}
		% \item[] \textbf{\#times} | list
			% % \begin{compactdesc}
				% % \item[Bible anchor verses] \phantom{.}
					% % \begin{compactdesc}
						% % \item[Romans 5:5] \gr{}
					% % \end{compactdesc}
				% % \item[Commentary] ---
			% % \end{compactdesc}
	% \end{compactdesc}
	
% % ----------------------
% \subsection{Bible Dictionaries} % *DICTIONARIES
% % ----------------------
	% \label{QUICKEN-BD}

% % ----------------------
% \subsection{Restoration Usage} % *RESTORATION
% % ----------------------
	% \label{QUICKEN-RESTORATION}

	% % -----
	% \subsubsection{Joseph Smith} % *JS
	% % -----
		% \label{QUICKEN-JS}
	
		% \begin{compactdesc}
			% \item[{\hyperref[KEY]{Date/Source}}]
		% \end{compactdesc}
		
	% % -----
	% \subsubsection{Temple Ordinances}
	% % -----
		% \label{QUICKEN-TEMPLE}
	
		% \begin{compactdesc}
			% \item[{\hyperref[KEY]{Date/Source}}]
		% \end{compactdesc}
	
	% % -----
	% \subsubsection{Brigham Young} % *BY
	% % -----
		% \label{QUICKEN-BY}
	
		% \begin{compactdesc}
			% \item[{\hyperref[KEY]{Date/Source}}]
		% \end{compactdesc}
	
% % ----------------------
% \subsection{Commentary and Studies} % *COMMENTARY *STUDIES
% % ----------------------
	% \label{QUICKEN-COMMENTARY}

	% % -----
	% \subsubsection{Key Points}
	% % -----
		% \label{QUICKEN-KEYS}
	
		% \begin{compactitem}
			% \item
		% \end{compactitem}
		
	% % -----
	% \subsubsection{Sample Study}
	% % -----
		% \label{QUICKEN-study}